\documentclass[11pt,a4paper]{article}
\input{../../sablony/pripravy/priprava.sty}
\usepackage{tikz,pgfplots}
\pgfplotsset{compat=1.18}
\nastavPripravu{1. ročník}{Kvadratické rovnice a nerovnice}{108}{Rozklad kvadratického trojčlenu}
\begin{document}
\blok{Rozklad kvadratického trojčlenu}
V minulé hodině jsme odvodili Viètovy vztahy. Nyní je spojíme se součinem kořenových činitelů:
\[a(x-x_1)(x-x_2)=ax^2-a(x_1+x_2)x+ax_1x_2=ax^2+bx+c.\]
Zpětná kontrola roznásobením je důležitější než mechanické hádání. 
Rozklad na součin a jeho využití při řešení rovnic; rozpoznání vhodné metody.
\begin{definiceBox}[title={Věta - Kořenový rozklad kvadratického trojčlenu}]
\textbf{Kořenový rozklad:} má-li $ax^2+bx+c$ kořeny $x_1,x_2$ (včetně případné násobnosti), platí $ax^2+bx+c=a(x-x_1)(x-x_2)$.
\end{definiceBox}
\blok{Metoda a návaznost na graf}
\begin{itemize}
\item Rozklad ověř roznásobením; není vždy jednodušší než diskriminant.
\item Je-li $D<0$, neexistuje kořenový rozklad na reálné lineární činitele.
\end{itemize}
\blok{Společně řešené příklady}
\begin{prikladReseniBox}[title={Příklad 1}]
$x^{2} - 7x + 10=0$
\tcblower
\[
\begin{aligned}
&\left(x - 5\right) \cdot \left(x - 2\right)=0\\[2pt]
x&=2\ \vee\ x=5\\[2pt]
K&=\{2;5\}
\end{aligned}
\]
\end{prikladReseniBox}
\begin{prikladReseniBox}[title={Příklad 2}]
$2x^{2} - 4x - 6=0$
\tcblower
\[
\begin{aligned}
&\left(x - 3\right) \cdot \left(2x + 2\right)=0\\[2pt]
x&=-1\ \vee\ x=3\\[2pt]
K&=\{-1;3\}
\end{aligned}
\]
\end{prikladReseniBox}
\begin{prikladReseniBox}[title={Příklad 3}]
$x^{2} + 2x + 5=0$
\tcblower
\[
\begin{aligned}
D&=4-4\cdot(1)\cdot(5)=-16\\[2pt]
D&<0\ \Rightarrow\ \text{žádný reálný kořen}\\[2pt]
K&=\varnothing
\end{aligned}
\]
\end{prikladReseniBox}
\blok{Samostatné procvičování}
\begin{zadaniReseniBox}[title={Samostatná práce 1}]
$x^{2} - x - 12=0$
\tcblower
\[
\begin{aligned}
&\left(x - 4\right) \cdot \left(x + 3\right)=0\\[2pt]
x&=-3\ \vee\ x=4\\[2pt]
K&=\{-3;4\}
\end{aligned}
\]
\end{zadaniReseniBox}
\begin{zadaniReseniBox}[title={Samostatná práce 2}]
$3x^{2} + 3x - 6=0$
\tcblower
\[
\begin{aligned}
&\left(x + 2\right) \cdot \left(3x - 3\right)=0\\[2pt]
x&=1\ \vee\ x=-2\\[2pt]
K&=\{-2;1\}
\end{aligned}
\]
\end{zadaniReseniBox}
\blok{Volitelná rezerva}
\poznamka{Studenti si vymyslí obdobný příklad se změnou jednoho koeficientu a vysvětlí, jak se změní počet kořenů či řešení.}
\blok{Rozšiřující příklady ze sbírky}
\poznamka{Jedna úloha je rozdíl čtverců, druhá se nad reálnými čísly na kořenové činitele nerozloží.}

\begin{zadaniReseniBox}[title={Rozšíření 1}]
Rozložte v $\mathbb R$ na kořenové činitele: $16x^{2} - 25$.
\tcblower
\[
\begin{aligned}P(x)&=16x^{2} - 25\\ &= 16x^{2} - 25\end{aligned}
\]
\end{zadaniReseniBox}

\begin{zadaniReseniBox}[title={Rozšíření 2}]
Rozložte v $\mathbb R$ na kořenové činitele: $x^{2} + 4$.
\tcblower
\[
\begin{aligned}P(x)&=x^{2} + 4\\ &= \text{nad }\mathbb R\text{ nelze rozložit na lineární činitele}\end{aligned}
\]
\end{zadaniReseniBox}

\blok{Domácí úkol}
\begin{zadaniReseniBox}[title={Domácí úkol}]
\begin{enumerate}[label=\alph*),itemsep=2pt]
\item $x^{2} - 8x + 7=0$
\item $2x^{2} - 4x - 16=0$
\item $3x^{2} + 15x=0$
\item $x^{2} + 5x + 4=0$
\item $- x^{2} + x + 6=0$
\item $2x^{2} + 4x + 9=0$
\item Rozložte v $\mathbb R$ na kořenové činitele: $2x^{2} + x - 1$.
\end{enumerate}
\tcblower
\begin{enumerate}[label=\alph*),itemsep=2pt]
\item
\[\begin{aligned}
&\left(x - 7\right) \cdot \left(x - 1\right)=0\\[2pt]
x&=1\ \vee\ x=7\\[2pt]
K&=\{1;7\}
\end{aligned}\]
\item
\[\begin{aligned}
&\left(x - 4\right) \cdot \left(2x + 4\right)=0\\[2pt]
x&=-2\ \vee\ x=4\\[2pt]
K&=\{-2;4\}
\end{aligned}\]
\item
\[\begin{aligned}
&3x \cdot \left(x + 5\right)=0\\[2pt]
x&=0\ \vee\ x=-5\\[2pt]
K&=\{-5;0\}
\end{aligned}\]
\item
\[\begin{aligned}
&\left(x + 1\right) \cdot \left(x + 4\right)=0\\[2pt]
x&=-4\ \vee\ x=-1\\[2pt]
K&=\{-4;-1\}
\end{aligned}\]
\item
\[\begin{aligned}
&\left(- x - 2\right) \cdot \left(x - 3\right)=0\\[2pt]
x&=-2\ \vee\ x=3\\[2pt]
K&=\{-2;3\}
\end{aligned}\]
\item
\[\begin{aligned}
D&=16-4\cdot(2)\cdot(9)=-56\\[2pt]
D&<0\ \Rightarrow\ \text{žádný reálný kořen}\\[2pt]
K&=\varnothing
\end{aligned}\]
\item \[ \begin{aligned}P(x)&=2x^{2} + x - 1\\ &= \left(x + 1\right) \cdot \left(2x - 1\right)\end{aligned}\]
\end{enumerate}
\end{zadaniReseniBox}
\poznamka{Orientační čas: úvod 5 min, společné příklady 18 min, samostatně 15 min, kontrola a shrnutí 7 min. Rezerva je dobrovolná. Číselné parametry modelových aplikací nejsou skutečně naměřená data.}
\end{document}
